\section{Discusión}
La contribución de un entorno declarativo es hacer explícitas las dependencias y ofrecer un punto de partida compartido. En formación en física, ese punto de partida permite dedicar una secuencia docente a decisiones científicas: selección de datos, unidades, normalización e interpretación. Enseñar Git y automatización añade un registro de decisiones que puede revisarse con el estudiante y reutilizarse en otras disciplinas.

El servicio remoto redistribuye trabajo y costos. Parte de la instalación pasa al equipo docente; cuentas, cuotas y conectividad siguen condicionando la sesión. Una imagen fijada por digest identifica el software distribuido, pero datos externos y servicios requieren su propia conservación. El mantenimiento debe combinar una versión estable para cada curso con una ruta explícita de actualización y prueba.

\paragraph{Alcance y límites} El trabajo es un ensayo técnico basado en documentación y diseño de arquitectura. No dispone de registros individuales de incorporación, exportaciones de consumo ni evaluación causal del aprendizaje de las escuelas citadas. Por ello no informa cifras históricas de velocidad, costos o incidencia de fallos. La comparación de plataformas describe responsabilidades y condiciones documentadas, no un ensayo de rendimiento; las plantillas proponen cómo obtener esa evidencia en futuras sesiones.

\section{Conclusiones}
La computación transversal se concreta en prácticas que sobreviven al cambio de plataforma: describir dependencias, versionar procedimientos, ejecutar comprobaciones y conservar resultados interpretables. La comparación identifica dónde recae la administración y qué debe prever una actividad de física para compartir un entorno de trabajo.

La evaluación futura debe reunir tres tipos de evidencia: ejecución limpia del análisis, consumo y soporte de la sesión, y comprensión del estudiante. Mantenerlos separados permite mejorar la infraestructura sin confundir rapidez de acceso con aprendizaje. Las plantillas de registro y el ejemplo editorial ofrecen instrumentos concretos para iniciar esa evaluación.

\section*{Vínculos profesionales}
Arturo Sánchez Pineda mantiene una relación profesional con INAIT SA y participa en iniciativas de formación y ciencia abierta citadas en este trabajo. Las afiliaciones de este manuscrito son Creative Commons Venezuela y Creative Commons Switzerland. \editorialtodo{La declaración de financiación específica se completará antes del envío.}
